Figure \ref{fig:jendl5-dnp-yield} shows the percentage of the total delayed neutron yield that comes from the six \glspl{DNP} with the largest yield contribution, where the yield from each \gls{DNP} is defined by Equation \eqref{eq:summation-tdny}.
There are 257 \glspl{DNP} present in JENDL-5, and they have a total delayed neutron yield of 0.0160.
This is consistent with the results in Section \ref{sec:iaea-jendl5}, as the only difference between these datasets is the emission probability and half-life data.

\begin{figure}
    \centering
    \includegraphics[width=0.5\linewidth]{images/jendl5/dnp_yield.png}
    \caption{The fraction of the delayed neutron yield from particular \glspl{DNP} from JENDL-5 nuclear data. The five \glspl{DNP} with the largest delayed neutron yield contributions approximately 52\% of the total delayed neutrons per fission event.}
    \label{fig:jendl5-dnp-yield}
\end{figure}


Table \ref{tab:jendl5-yld-hl} shows how the total delayed neutron yield $\nu_d$ and average half-life $\bar{\tau}$ compares between the value from the $I$ \glspl{DNP} and the six \gls{DNP} groups.
The results agree well with the data in Section \ref{sec:iaea-jendl5}.

\begin{table}[]
    \centering
    \caption{JENDL-5 nuclear data delayed neutron yield and average half-life values.}
    \begin{tabular}{lrl}
        \hline
        Parameter & Value & Units\\
        \hline
        $\nu_d(I)$ & $0.0160 \pm 0.0009$ &$ [-]$\\
        $\nu_d(K)$ & $0.0160 \pm 0.0008$ & $[-]$\\
        $\Delta \nu_d$ & $0.09 \pm 122.31$ & $[pcm]$\\
        \hline
        $\bar{\tau} (I)$ & $9.1 \pm 0.5$ & $[s]$\\
        $\bar{\tau} (K)$ & $9.1 \pm 0.5$ & $[s]$\\
        $\Delta \bar{\tau}$ & $0.0 \pm 0.6$ & $[s]$\\ % 0.00591310092491959+/-0.6488314744438124 s
        \hline
    \end{tabular}
    \label{tab:jendl5-yld-hl}
\end{table}

Table \ref{tab:jendl5-PCCi} shows the ten largest $PCC_i$ terms, which reveals which \glspl{DNP} have the greatest impact on the \gls{DNP} group parameters $\nu_{d,k}$ and $\tau_k$.
The same \glspl{DNP} are in Table \ref{tab:iaea-jendl5-PCCi}, though the exact values and order is slightly different due to half-life and emission probability differences.


\begin{table}[]
    \centering
    \caption{JENDL-5 nuclear data ten largest \glspl{PCC}.}
    \begin{tabular}{ll}
        \hline
       Nuclide & $PCC_{i}$ \\
       \hline
       $^{87}$Br & 3.788538 \\
       $^{88}$Br & 3.405082 \\
       $^{137}$I & 3.373592 \\
       $^{96}$Rb & 3.230474 \\
       $^{97}$Rb & 2.406196 \\
       $^{138}$I & 2.209843 \\
       $^{94}$Rb & 2.202741 \\
       $^{89}$Br & 2.123136 \\
       $^{91}$Br & 1.779122 \\
       $^{90}$Br & 1.733836 \\
        \hline
    \end{tabular}
    \label{tab:jendl5-PCCi}
\end{table}

Table \ref{tab:jendl5-Ui} shows the ten largest $U_i$ terms, which corresponds to which \glspl{DNP} have the largest uncertainties relative to their impact on the \gls{DNP} group parameters $\nu_{d,k}$ and $\tau_k$.
Four out of ten of these \glspl{DNP} are also in Table \ref{tab:iaea-jendl5-Ui}, though there are differences due to the emission probability and half-life differences.
One notable difference is $^{102}$Sr, which has a $U_i$ value of $3.63 \times 10^{13}$.
This value is significantly larger than any other value from any other dataset because it has an emission probability of $(6.245 \times 10^{-17}) \pm 0.02$.
This means that the relative uncertainty in Equation \eqref{eq:Uiv} works out to be $3.2 \times 10^{-14}$, which is exceedingly large.
To continue conducting the analysis, the contribution of $^{102}$Sr is ignored, treating the emission probability as 0 for this analysis.

\begin{table}[]
    \centering
    \caption{JENDL-5 nuclear data ten largest uncertainty-weighted \glspl{PCC}.}
    \begin{tabular}{ll}
        \hline
       Nuclide & $U_{i}$ \\
       \hline
       $^{102}$Sr & -\\
       $^{142}$Sb & 0.283440 \\
       $^{99}$Kr & 0.279344 \\
       $^{105}$Zr & 0.270256 \\
       $^{116}$Rh & 0.269286 \\
       $^{103}$Zr & 0.267184 \\
       $^{124}$Ag & 0.228787 \\
       $^{68}$Co & 0.224881 \\
       $^{134}$Sn & 0.207609 \\
       $^{66}$Mn & 0.199226 \\
        \hline
    \end{tabular}
    \label{tab:jendl5-Ui}
\end{table}

Figure \ref{fig:jendl5-bar-uiv} gives a more detailed view on which \gls{DNP} values specifically lead to the large $U_i$ values.
This offers slightly more insight than Table \ref{tab:jendl5-Ui}, which sums over the various \gls{DNP} values.
These uncertainty-weighted sensitivities are closer to the ideal configuration: values are close to each other, and the uncertainties in half-lives, emission probability, and concentration each contribute some amount.
There are no clear outliers in the data, though, similar to the results in Section \ref{sec:iaea-jendl5}, the largest uncertainty-weighted sensitivity contributions come from emission probabilities, indicating a potential area of improvement.


\begin{figure}
    \centering
    \includegraphics[width=0.5\linewidth]{images/jendl5/pcc-bar.png}
    \caption{The ten largest uncertainty-weighted \glspl{PCC} $U_{i,v}$ sorted by \gls{DNP} and associated \gls{DNP} value from ENDF/B-VII.1 nuclear data.}
    \label{fig:jendl5-bar-uiv}
\end{figure}


Figure \ref{fig:pcc-charts-jendl5} shows the $PCC_i$ and $U_i$ plotted on the chart of the nuclides.
These plots do not limit to only the top ten largest values, but instead show how all of the values compare among each other.
The values of $U_i$ are all smaller than those of $PCC_i$, because all of the relative uncertainties are less than 1.
There are similar patterns as discussed in Section \ref{sec:iaea-jendl5}, with horizontal peaks in the $PCC_i$ for bromine, rubidium, and iodine.

\begin{figure}[!ht]
\centering
\subfloat[The $PCC_i$, or total \gls{PCC}, calculated for each \gls{DNP}.]{
  \includegraphics[width=0.45\linewidth]{images/jendl5/chart_PCC.png}
}
\hspace{0.05\linewidth}
\subfloat[The $U_i$, or total uncertainty-weighted \gls{PCC}, for each \gls{DNP}.]{
  \includegraphics[width=0.45\linewidth]{images/jendl5/chart_PCC_uncertainty.png}
}
\caption{Chart of nuclide plots showing the \glspl{PCC} and uncertainty-weighted \glspl{PCC} using JENDL-5 nuclear data.}
\label{fig:pcc-charts-jendl5}
\end{figure}

\FloatBarrier


